\paragraph{Cited theorem (relative consistency of the negation of CH).}
If \(\ZFC\) is consistent, then \(\ZFC+\neg\mathsf{CH}\) is consistent
\cite[Chapter~14, pp.~201--224]{Jech2003}.  The import is a family of
\(\ZFC\)-theorem schemes about the definable class \(V^{B}\) of
\(B\)-valued names, where \(B\) is the Boolean completion of the poset
\(\operatorname{Fn}(\aleph_2\times\omega,2,{<}\omega)\) of finite partial
functions from \(\aleph_2\times\omega\) to \(2\), and
\(\lVert\sigma\rVert_{B}\) denotes the defined \(\ZFC\)-term assigning to
a membership-language sentence \(\sigma\) with name parameters its Boolean
truth value in \(B\); both notations are scoped to this import.  The
cited schemes are:
\begin{enumerate}
  \item \(\lVert\sigma\rVert_{B}=1\) for each axiom \(\sigma\) of
        \(\ZFC\), including every schema instance;
  \item Boolean soundness: if
        \(\set{\sigma_1,\ldots,\sigma_k}\derives\psi\), then
        \(\ZFC\derives
          \lVert\sigma_1\rVert_{B}\wedge\cdots\wedge
          \lVert\sigma_k\rVert_{B}\leq\lVert\psi\rVert_{B}\),
        where \(\lVert\bottom\rVert_{B}=0\) and \(\ZFC\) proves \(B\)
        nontrivial, that is, \(0\ne1\) in \(B\);
  \item \(B\) satisfies the countable chain condition, and
        chain-condition forcing preserves cardinals:
        \(\lVert\check\aleph_1=\aleph_1\land
           \check\aleph_2=\aleph_2\rVert_{B}=1\); and
  \item the distinctness density argument of the one-Cohen-real example,
        computed per pair of \(\aleph_2\)-indexed sections of the adjoined
        function, gives
        \(\lVert 2^{\aleph_0}\geq\check\aleph_2\rVert_{B}=1\), hence
        \(\lVert\neg\mathsf{CH}\rVert_{B}=1\).
\end{enumerate}
A derivation of \(\bottom\) from \(\ZFC\cup\set{\neg\mathsf{CH}}\) would
therefore give \(\ZFC\derives 1=0\) in \(B\), contradicting
nontriviality.  This Boolean-valued route is the discharge demanded by
the warning on countable transitive structures below; it hypothesizes no
countable transitive model.  The source states these schemes for its own
calculus; consistency transfers to the fixed calculus \(\ND_{\FOL}\)
because soundness (\cref{thm:fol-soundness}) and strong completeness
(\cref{thm:fol-completeness}) make derivability calculus-invariant.

